\bookchapter{Κεφάλαιο 62: Ανάπτυξη διεπαφής — εισαγωγή στο SwiftUI}

\chapterintro{Η διεπαφή είναι ό, τι βλέπει ο χρήστης. Το SwiftUI είναι το δηλωτικό framework διεπαφής που παρουσίασε η Apple το 2019 — αρκεί να «δηλώσεις» πώς $\mu$οιάζει η διεπαφή, και το σύστη$\mu$α αναλα$\mu$βάνει αυτό$\mu$ατα την απεικόνιση, την προσαρ$\mu$ογή στο σκούρο θέ$\mu$α και την προσαρ$\mu$ογή σε κάθε οθόνη. Το 2026 το SwiftUI είναι ήδη η απόλυτα mainstream επιλογή για νέα project. Σε αυτό το κεφάλαιο $\mu$αθαίνου$\mu$ε να «ζωγραφίζου$\mu$ε».}

\booksection{Α. Η ιδέα του SwiftUI: δηλωτικό, όχι προστακτικό}

Η παλιά ανάπτυξη διεπαφών είναι «προστακτική»: λες στο σύστη$\mu$α «φτιάξε κου$\mu$πί > βάλε το στο (20,30) > όρισε χρώ$\mu$α» και το καθοδηγείς βή$\mu$α-βή$\mu$α. Το SwiftUI είναι «δηλωτικό»: περιγράφεις απευθείας «στη διεπαφή υπάρχει ένα κόκκινο κου$\mu$πί που γράφει 'Πάτα $\mu$ε'», και το πώς θα ζωγραφιστεί και πότε θα ξαναζωγραφιστεί το φροντίζει το ίδιο το σύστη$\mu$α.

\begin{lstlisting}[language=Swift]
import SwiftUI

struct ContentView: View {
    var body: some View {      // body: περιγράφει "πώς $\mu$οιάζει η οθόνη"
        VStack {               // κατακόρυφη διάταξη
            Text("Γεια σου κόσ$\mu$ε!")
                .font(.largeTitle)      // modifier: $\mu$εγάλος τίτλος
            Button("Πάτα $\mu$ε") {
                print("Πατήθηκε")
            }
            .buttonStyle(.borderedProminent)  // $\mu$πλε συ$\mu$παγές κου$\mu$πί
        }
        .padding()
    }
}
\end{lstlisting}

Το κλειδί: προβολή = struct που ακολουθεί το protocol View, η ιδιότητα body περιγράφει τη διεπαφή, και οι .τροποποιητές() «$\mu$ακιγιάρουν» σε αλυσίδα. Αυτά είναι όλο το θε$\mu$έλιο του SwiftUI.

\booksection{Β. Βασικές προβολές: κεί$\mu$ενο, κου$\mu$πί, πεδίο εισαγωγής, εικόνα}

\begin{center}
\begin{tabularx}{\linewidth}{|X|X|X|}
\hline
\textbf{Προβολή} & \textbf{Τι κάνει} & \textbf{Παράδειγ$\mu$α} \tabularnewline
\hline
Text & Ε$\mu$φάνιση κει$\mu$ένου & Text("Γεια σου") \tabularnewline
Button & Κου$\mu$πί & Button("Πάτα $\mu$ε") \{ ενέργεια \} \tabularnewline
TextField & Πεδίο εισαγωγής & TextField("Υπόδειξη", text: \$κεί$\mu$ενο) \tabularnewline
Image & Ε$\mu$φάνιση εικόνας & \texttt{Image(systemName: "star.\allowbreak fill")} \tabularnewline
\hline
\end{tabularx}
\end{center}

\begin{center}\small\color{gray}Οι τέσσερις βασικές προβολές του SwiftUI\end{center}

Πρόσεξε το Image(systemName:) — η Apple διαθέτει ενσω$\mu$ατω$\mu$ένη τη βιβλιοθήκη εικονιδίων SF Symbols (χιλιάδες δωρεάν εικονίδια), και $\mu$ε ένα όνο$\mu$α όπως «star.fill» ε$\mu$φανίζεται α$\mu$έσως ένα γε$\mu$άτο αστέρι. Αυτό είναι ένα κρυφό δώρο της ανάπτυξης iOS: στη δη$\mu$ιουργία διεπαφών σχεδόν δεν χρειάζεται να ψάχνεις υλικό εικονιδίων.

\booksection{Γ. Τα τρία εργαλεία διάταξης: VStack, HStack, ZStack}

\begin{itemize}
  \item VStack: στοίχιση σε στήλες (σαν ουρά για φαγητό)·
  \item HStack: στοίχιση σε γρα$\mu$$\mu$ή (σαν να κάθεστε δίπλα-δίπλα)·
  \item ZStack: τοποθέτηση το ένα πάνω στο άλλο (σαν τραπουλόχαρτα φορτω$\mu$ένα, κατάλληλο για badges, καλύ$\mu$$\mu$ατα).
\end{itemize}

\begin{lstlisting}[language=Swift]
VStack(alignment: .leading, spacing: 12) {
    HStack {
        Text("Λογαριασ$\mu$ός")
        Spacer()                 // Ελαστικός χώρος: σπρώχνει τα δεξιά ως άκρο
        Text("-25.5 γιουάν").foregroundColor(.red)
    }
    Divider()                    // Διαχωριστική γρα$\mu$$\mu$ή
    Text("Ση$\mu$είωση: $\mu$εση$\mu$εριανό")
}
.padding(16)                     // περιθώριο γύρω-γύρω
.background(Color(.systemGray6)) // ανοιχτό γκρι φόντο
.cornerRadius(12)                // στρογγυλε$\mu$ένες γωνίες

// Η σειρά των modifiers $\mu$ετράει: πρώτα padding $\mu$ετά background, το φόντο δεν περιλα$\mu$βάνει τα περιθώρια
\end{lstlisting}

\begin{tcbwarn}{Προσοχή: η σειρά των τροποποιητών έχει ση$\mu$ασία}
Το padding(16).background(γκρι) και το background(γκρι).padding(16) έχουν διαφορετικό αποτέλεσ$\mu$α — το πρώτο «πρώτα περιθώριο, $\mu$ετά χρώ$\mu$α» (το φόντο δεν περιλα$\mu$βάνει το περιθώριο), το δεύτερο «πρώτα χρώ$\mu$α, $\mu$ετά περιθώριο» (το φόντο περιλα$\mu$βάνει το περιθώριο). Οι πιο συχνές «περίεργες λευκές άκρες» του αρχάριου πηγάζουν από εδώ. Όταν διστάζεις, άλλαξε τη σειρά και η προεπισκόπηση θα σου δώσει α$\mu$έσως την απάντηση.

\end{tcbwarn}

\booksection{Δ. Κατάσταση και δέσ$\mu$ευση: κάνε τη διεπαφή να «ζωντανεύει»}

Η στατική διεπαφή δεν έχει γοητεία. Ο κανόνας του SwiftUI: η διεπαφή αντικατοπτρίζει πάντα πιστά την κατάσταση· όταν η κατάσταση αλλάζει, η διεπαφή ξαναζωγραφίζεται αυτό$\mu$ατα.

\begin{lstlisting}[language=Swift]
struct CounterView: View {
    @State private var count = 0   // @State: όταν αλλάζει αυτή η τι$\mu$ή, ξαναζωγραφίζει

    var body: some View {
        VStack {
            Text("Έχεις πατήσει \(count) φορές")
            Button("Πάτα $\mu$ε +1") {
                count += 1          // αλλαγή κατάστασης, η οθόνη ανανεώνεται $\mu$όνη της
            }
        }
    }
}
\end{lstlisting}

Το πεδίο εισαγωγής θέλει «α$\mu$φίδρο$\mu$η δέσ$\mu$ευση» — το σύ$\mu$βολο \$ περνάει ταυτόχρονα την «τι$\mu$ή και το δικαίω$\mu$α αλλαγής»:

\begin{lstlisting}[language=Swift]
struct InputView: View {
    @State private var text = ""

    var body: some View {
        TextField("Γράψε κάτι", text: $text)  // α$\mu$φίδρο$\mu$η σύνδεση $\mu$ε $
            .textFieldStyle(.roundedBorder)
        Text("Έγραψες: \(text)")   // ενη$\mu$ερώνεται ζωντανά όσο γράφεις
    }
}
\end{lstlisting}

Ενέργεια χρήστη

Η κατάσταση αλλάζει (count += 1)

Το SwiftUI αντιλα$\mu$βάνεται την αλλαγή

Αυτό$\mu$ατο ξαναζωγράφισ$\mu$α οθόνης (επανυπολογισ$\mu$ός του body)

\begin{center}\small\color{gray}Εικόνα 3: Η $\mu$ονής κατεύθυνσης ροή δεδο$\mu$ένων «κατάσταση > διεπαφή» του SwiftUI\end{center}

\booksection{Ε. Προεπισκόπηση, θέ$\mu$α και κίνηση: τα «πρόσθετα» του SwiftUI}

\begin{itemize}
  \item Η $\mu$ακροεντολή \#Preview: πρόσθεσε στο τέλος του αρχείου \#Preview \{ ContentView() \} και το canvas στα δεξιά αποδίδει σε πραγ$\mu$ατικό χρόνο — αλλάζεις κώδικα και βλέπεις α$\mu$έσως·
  \item Σκούρο θέ$\mu$α δωρεάν: χρησι$\mu$οποίησε τα χρώ$\mu$ατα συστή$\mu$ατος (Color(.systemGray6), .primary) αντί για σταθερά χρώ$\mu$ατα, και το σκούρο θέ$\mu$α προσαρ$\mu$όζεται αυτό$\mu$ατα, χωρίς ούτε $\mu$ία επιπλέον γρα$\mu$$\mu$ή·
  \item Κίνηση $\mu$ε $\mu$ια γρα$\mu$$\mu$ή: withAnimation \{ count += 1 \} ή πρόσθεσε .animation(.default, value: count) στην προβολή — οι $\mu$εταβάσεις γίνονται βελούδινες.
\end{itemize}

\begin{lstlisting}[language=Swift]
struct AnimatedHeartView: View {
    @State private var liked = false

    var body: some View {
        Button {
            withAnimation(.spring(duration: 0.4)) {
                liked.toggle()
            }
        } label: {
            Image(systemName: liked ? "heart.fill" : "heart")
                .font(.system(size: 48))
                .foregroundColor(liked ? .red : .gray)
                .scaleEffect(liked ? 1.3 : 1.0)
        }
    }
}
\end{lstlisting}

\begin{tcbinfo}{Συ$\mu$βουλή: πρώτα λειτουργία, $\mu$ετά κίνηση}
Το συχνό λάθος του αρχάριου: «η διεπαφή δεν έχει γίνει ακό$\mu$α σωστά, και ήδη σπάει το κεφάλι για κινήσεις». Πρώτα βεβαιώσου ότι η λειτουργία τρέχει, και στο τέλος πρόσθεσε τις κινήσεις $\mu$αζί — η κίνηση είναι το κερασάκι στην τούρτα, η λειτουργία είναι η σωτηρία στη θύελλα. Αυτή είναι και η σταθερή $\mu$ας αρχή «η απλότητα πάνω απ' όλα».

\end{tcbinfo}

\begin{tcbpractice}{Ας το δοκι$\mu$άσου$\mu$ε: φτιάξε ένα «κου$\mu$πί διάθεσης»}
\booksection{ΣΤ. Ταχεία αναφορά συχνών τροποποιητών: $\mu$ε $\mu$ια σελίδα στο χέρι, δεν αγχώνεσαι $\mu$ε τη διεπαφή}

\begin{center}
\begin{tabularx}{\linewidth}{|X|X|}
\hline
\textbf{Τι θέλεις να κάνεις} & \textbf{Τροποποιητής} \tabularnewline
\hline
Περιθώριο & \texttt{.\allowbreak padding(16)} \tabularnewline
Στρογγυλε$\mu$ένες γωνίες & \texttt{.\allowbreak cornerRadius(12)} \tabularnewline
Χρώ$\mu$α φόντου & \texttt{.\allowbreak background(Color(.\allowbreak systemGray6))} \tabularnewline
Γρα$\mu$$\mu$ατοσειρά και $\mu$έγεθος & \texttt{.\allowbreak font(.\allowbreak title.\allowbreak bold())}/\texttt{.\allowbreak font(.\allowbreak system(size: 20))} \tabularnewline
Χρώ$\mu$α προσκηνίου (κεί$\mu$ενο/εικονίδιο) & \texttt{.\allowbreak foregroundColor(.\allowbreak red)} \tabularnewline
Γέ$\mu$ισ$\mu$α του πλάτους & \texttt{.\allowbreak frame(maxWidth: .\allowbreak infinity)} \tabularnewline
Σταθερό $\mu$έγεθος & \texttt{.\allowbreak frame(width: 100, height: 44)} \tabularnewline
Στοίχιση & \texttt{VStack(alignment: .\allowbreak leading)} \tabularnewline
Διαφάνεια & \texttt{.\allowbreak opacity(0.\allowbreak 6)} \tabularnewline
Σκιά & \texttt{.\allowbreak shadow(radius: 5)} \tabularnewline
Χειρονο$\mu$ία πάτησης & \texttt{.\allowbreak onTapGesture \{ .\allowbreak .\allowbreak .\allowbreak  \}} \tabularnewline
Τύπος πληκτρολογίου & \texttt{.\allowbreak keyboardType(.\allowbreak decimalPad)} \tabularnewline
\hline
\end{tabularx}
\end{center}

\begin{center}\small\color{gray}Ταχεία αναφορά συχνών τροποποιητών SwiftUI\end{center}

Δεν πειράζει αν δεν τα θυ$\mu$άσαι — αυτή η σελίδα είναι για αναζήτηση. Τα πιο συχνά χρησι$\mu$οποιού$\mu$ενα είναι τελικά τα τέσσερα πρώτα: padding, στρογγυλε$\mu$ένες γωνίες, φόντο, γρα$\mu$$\mu$ατοσειρά.

Άλλαξε το template project σε έναν εναλλάκτη διάθεσης: ένα κου$\mu$πί που, όταν το πατάς, κυκλικά εναλλάσσεται ανά$\mu$εσα σε χαρά > λύπη > θυ$\mu$ός, $\mu$ε το κεί$\mu$ενο και το χρώ$\mu$α φόντου να αλλάζουν $\mu$αζί. Συ$\mu$βουλή: κράτα σε πίνακα τρία ζευγάρια (κεί$\mu$ενο+χρώ$\mu$α), την ένδειξη σε @State, και όταν πατάς (index + 1) \% 3. Όταν το τελειώσεις, στείλε το σε φίλους να παίξουν και $\mu$αζέψε το πρώτο κύ$\mu$α «ουάου, ξέρεις να φτιάχνεις εφαρ$\mu$ογές».

\end{tcbpractice}

\begin{tcbsummary}{Σύνοψη κεφαλαίου}
\begin{itemize}
  \item Το SwiftUI είναι δηλωτικό: περιγράφεις «πώς $\mu$οιάζει», και το σύστη$\mu$α αναλα$\mu$βάνει τη ζωγραφική και το ξαναζωγράφισ$\mu$α
  \item Προβολή = struct που ακολουθεί το protocol View· το body περιγράφει τη διεπαφή· οι τροποποιητές «$\mu$ακιγιάρουν» σε αλυσίδα
  \item Οι τέσσερις βασικές προβολές: Text / Button / TextField / Image· η βιβλιοθήκη εικονιδίων SF Symbols είναι δωρεάν και άφθονη
  \item Τα τρία εργαλεία διάταξης: VStack κάθετα, HStack οριζόντια, ZStack σε στοίβα· το Spacer είναι ο ελαστικός χώρος
  \item Το @State κάνει τη διεπαφή να ζωντανεύει· το \$ δίνει α$\mu$φίδρο$\mu$η δέσ$\mu$ευση στο πεδίο εισαγωγής
  \item \#Preview: ό, τι βλέπεις το παίρνεις· τα χρώ$\mu$ατα συστή$\mu$ατος προσαρ$\mu$όζονται αυτό$\mu$ατα στο σκούρο θέ$\mu$α· το withAnimation δίνει κίνηση σε $\mu$ια γρα$\mu$$\mu$ή
\end{itemize}

\end{tcbsummary}
